% BHU PYQ -- Manually transcribed & error-corrected
% Paper: GLB-501: Physical and Structural Geology
% Exam: B.Sc. (Hons.) Semester V Examination, 2016-17

\documentclass[11pt,a4paper]{article}
\usepackage[a4paper,top=2.5cm,bottom=2.5cm,left=2.8cm,right=2.8cm,headheight=14pt]{geometry}
\usepackage{amsmath,amssymb}
\usepackage[T1]{fontenc}
\usepackage[utf8]{inputenc}
\usepackage{lmodern,microtype,fancyhdr,enumitem,hyperref,lastpage,parskip}

\hypersetup{colorlinks=true,urlcolor=black,linkcolor=black,
  pdftitle={GLB-501 Physical and Structural Geology -- BHU B.Sc. Sem V 2016-17}}
\pagestyle{fancy}\fancyhf{}
\renewcommand{\headrulewidth}{0.4pt}\renewcommand{\footrulewidth}{0.4pt}
\fancyhead[L]{\small   \textit{Banaras Hindu University}}
\fancyhead[R]{\small   \textit{GLB-501: Physical and Structural Geology}}
\fancyfoot[C]{\small Page  \thepage\ of \pageref{LastPage}}


\newlist{parts}{enumerate}{2}
\setlist[parts,1]{label=(\alph*),leftmargin=*,itemsep=5pt,topsep=3pt}

\newlist{romanparts}{enumerate}{2}
\setlist[romanparts,1]{label=(\roman*),leftmargin=*,itemsep=5pt,topsep=3pt}

\newcommand{\pts}[1]{\hfill   \textit{[#1]}}


\begin{document}

\begin{center}
  {\large\bfseries BANARAS HINDU UNIVERSITY}\\[2pt]
  {\normalsize B.Sc. (Hons.) --- Semester V Examination, 2016-17}\\[6pt]
  \rule{0.8\linewidth}{0.5pt}\\[6pt]
  {\large\bfseries Paper No. GLB-501: Physical and Structural Geology}\\[4pt]
  {\normalsize Time: 3 Hours \hfill Maximum Marks: 70}\\[2pt]
  {\small Ref. No.: BS/Sem V/175}
\end{center}
\rule{\linewidth}{0.4pt}
\smallskip



\noindent   \textit{Instructions: Attempt any   \textbf{five} questions in all, at least   \textbf{two} questions each from Section A and B. Section C is compulsory. All questions carry equal marks (14 marks each).}
\bigskip

\bigskip


\noindent {\large\bfseries SECTION --- A}
\rule{\linewidth}{0.4pt}\smallskip




\noindent   \textbf{Question 1.} Describe the classification of folds based on dip-isogons. \pts{14}

\medskip


\noindent   \textbf{Question 2.} Describe the erosional and depositional features formed by rivers. \pts{14}

\medskip


\noindent   \textbf{Question 3.} Discuss with suitable examples: \pts{$2  \times 7 = 14$}

\begin{parts}
  \item Types of sand dunes
  \item Mitigation of landslides
\end{parts}

\medskip


\noindent   \textbf{Question 4.} What is an unconformity? Describe the different types of unconformity and explain how you would recognize an unconformity in the field. \pts{14}

\bigskip


\noindent {\large\bfseries SECTION --- B}
\rule{\linewidth}{0.4pt}\smallskip




\noindent   \textbf{Question 5.} Write detailed notes on the following: \pts{$2  \times 7 = 14$}

\begin{parts}
  \item Effect of faulting on folded outcrop
  \item Planar and linear structures in deformed rocks
\end{parts}

\medskip


\noindent   \textbf{Question 6.} Write short notes on the following: \pts{$4  \times 3.5 = 14$}

\begin{parts}
  \item Pratt's concept of Isostasy
  \item Erosional geomorphic features of the ocean
  \item Divergent plate boundary
  \item Plateau
\end{parts}

\medskip


\noindent   \textbf{Question 7.} Discuss the following: \pts{$2  \times 7 = 14$}

\begin{parts}
  \item Endogenic processes
  \item Joints and their geometrical classification
\end{parts}

\medskip


\noindent   \textbf{Question 8.} Differentiate between the following: \pts{$4  \times 3.5 = 14$}

\begin{parts}
  \item Window and Klippe
  \item Horst and Graben
  \item Tension joint and shear joint
  \item Fault and unconformity
\end{parts}

\bigskip


\noindent {\large\bfseries SECTION --- C (Compulsory)}
\rule{\linewidth}{0.4pt}\smallskip




\noindent   \textbf{Question 9.} Describe each of the following in two to three sentences. Give suitable diagrams where applicable: \pts{$7  \times 2 = 14$}

\begin{romanparts}
  \item Geomagnetic time scale
  \item Difference between till and moraine
  \item Angle of repose
  \item Allochthonous tectonic unit
  \item Flower structure
  \item Similar fold and Parallel fold
  \item Antiformal Syncline
\end{romanparts}

\vfill\rule{\linewidth}{0.4pt}

\begin{center}\small
\href{https://bhu-exam.vercel.app/}{Click here to visit our website}\\[4pt]
\href{https://me-aryan.vercel.app/}{Click here to visit our contributor}\\[4pt]
\href{https://quashan-me.vercel.app/}{Click here to visit our contributor}
\end{center}
\end{document}
